\chapter{Synthèse de la correspondance}

\section{Tableau récapitulatif de couverture des blocs de compétences RNCP38152}

\begin{landscape}

\begin{longtable}{|p{4cm}|p{9cm}|p{11cm}|}
\caption{Synthèse de la correspondance parcours - RNCP38152} \\
\hline
\textbf{Bloc RNCP} & \textbf{Expériences principales} & \textbf{Preuves tangibles} \\
\hline
\endfirsthead
\hline
\textbf{Bloc RNCP} & \textbf{Expériences principales} & \textbf{Preuves tangibles} \\
\hline
\endhead
\hline
\endfoot

\textbf{BC01 - Usages numériques} &
20 ans d'expertise technique, développement d'EIAH, conception d'outils sur mesure (moteur de jeu, outil d'évaluation) &
Applications sur la Forge, Code source du moteur de jeu, Plateforme "Notation à l'Enveloppe" \\
\hline

\textbf{BC02 - Savoirs spécialisés} &
Double compétence expert technique / ingénieur pédagogique, innovation méthodologique &
Diplômes techniques, méthode "Notation à l'Enveloppe", productions didactiques \\
\hline

\textbf{BC03 - Communication} &
Formation de formateurs, documentation technique et pédagogique, vulgarisation &
Supports de formation, documentation Forge, site web éducatif \\
\hline

\textbf{BC04 - Transformation pro} &
Pilotage de projets complexes (industrie \& éducation), gestion d'équipes pluridisciplinaires &
Projet Jeu Vidéo (5 enseignants, 40 élèves), management d'équipes techniques (12 ans) \\
\hline

\textbf{BC05 - Ingénierie formation} &
Conception de dispositifs complets (SQL, Jeu), analyse des besoins, scénarisation &
Scénarios pédagogiques, grilles d'évaluation, plans de formation \\
\hline

\textbf{BC06 - Écosystème \& Coopération} &
Travail collaboratif en réseau (Forge, NSI), projets interdisciplinaires &
Contributions open source, projet inter-matières, partenariats académiques \\
\hline

\textbf{BC07 - Production ressources} &
Création et diffusion de ressources numériques innovantes (REL), outils d'évaluation &
10+ applications libres, tutoriels vidéo, fiches pédagogiques \\
\hline

\textbf{BC08 - Analyse réflexive} &
Accompagnement de pairs, auto-analyse systématique des dispositifs, démarche qualité &
Bilans d'expériences (Partie 2), ajustements des scénarios, mentorat informel \\
\hline

\textbf{BC09 - Développement pro} &
Veille scientifique (IA, didactique), formation continue, recherche-action &
Projet de thèse, bibliographie, participation à des groupes de travail \\
\hline

\end{longtable}

\end{landscape}

\section{Conclusion : Une expertise en ingénierie de formation}

L'analyse de la correspondance entre mon parcours et le référentiel RNCP38152 met en évidence une couverture étendue des 9 blocs de compétences attendus pour le Master MEEF PIF. Cette adéquation repose sur trois piliers complémentaires :

\subsection{Expertise technologique au service de la formation}

\begin{itemize}[leftmargin=*]
    \item 20 ans d'expérience technique de haut niveau réinvestis dans la conception d'outils pour la formation (BC01, BC07).
    \item Une capacité à anticiper les évolutions numériques (IA) pour adapter les ingénieries de formation.
\end{itemize}

\subsection{Compétences en ingénierie pédagogique}

\begin{itemize}[leftmargin=*]
    \item Conception et pilotage de dispositifs de formation complexes et innovants (BC04, BC05).
    \item Une approche systémique de la formation, intégrant l'analyse des besoins, la conception, l'animation et l'évaluation.
\end{itemize}

\subsection{Posture de praticien-chercheur et formateur de formateurs}

\begin{itemize}[leftmargin=*]
    \item Une démarche réflexive constante et partagée (BC08, BC09).
    \item Une contribution active à la communauté éducative par la production de savoirs et de ressources (BC02, BC06).
\end{itemize}

\vspace{1cm}

Cette triple légitimité constitue un socle solide pour valider le Master MEEF PIF et poursuivre vers une carrière d'ingénieur de formation ou d'enseignant-chercheur.
